\documentclass[10pt]{article}
\usepackage[letterpaper,margin=0.9in]{geometry}
\usepackage{amsmath,amssymb,amsthm,booktabs,microtype}
\usepackage[colorlinks=true,linkcolor=blue,citecolor=blue,urlcolor=blue]{hyperref}
\newtheorem{proposition}{Proposition}
\newcommand{\arcosh}{\operatorname{arcosh}}
\newcommand{\esssup}{\operatorname*{ess\,sup}}
\newcommand{\R}{\mathbb R}
\title{A finite-flux speed limit for Page-like entropy excursions}
\author{}
\date{27 October 2026}
\begin{document}
\maketitle
\vspace{-1.7em}
\begin{abstract}
How rapidly can a complete entropy excursion occur if the outgoing null-energy
flux is bounded in magnitude? For two-dimensional conformal vacua, we give the
sharp duration bound at fixed entropy height and independent positive and
negative flux caps. Its optimizer has two positive-flux intervals enclosing one
negative-flux interval, with no distributional energy pulses. The equal-cap
problem is an application of the known minimum-time quantum-piston control
problem; we make that connection explicit and give a direct proof covering
arbitrary additional turnovers. The bound is sharp as an infimum over smooth
conformal vacua with the same height and caps. We also derive a necessary lower
bound on peak quantum-null-energy-condition nonsaturation from an entropy
excursion faster than the conformal-vacuum limit. The result concerns a specified
null entropy and affine-time normalization, rather than a general bound on
black-hole evaporation.
\end{abstract}

\section{Question and setting}

In two-dimensional conformal vacua, the outgoing energy flux and renormalized
radiation entropy obey an exact differential relation \cite{BS}. This permits
entropy histories to constrain energy emission without specifying a bulk
evaporation model. Integrated-energy costs of complete Page-like excursions,
including their tradeoffs, have recently been optimized \cite{KD}; earlier
work discussed the energy--entropy relation in evaporation \cite{AP}. Here the
constraint is instead on the \emph{instantaneous} positive and negative flux.
An optimum involving delta-function pulses cannot satisfy such a constraint.

We develop the conformal-field-theory application, state the unequal-cap
extension, and derive a nonsaturation diagnostic. The equal-cap control problem was
solved by Stefanatos for a quantum piston \cite{Stefanatos}. Hyperbolic and
trigonometric matching of square energy pulses also occurs in quantum-interest
analysis \cite{TW}. We do not claim a new optimal-control principle or a new
derivation of the conformal anomaly.

Set $\hbar=1$ and let $c>0$ denote the central charge. For an asymptotically
affine ray-tracing map $p(u)$, with $p'>0$ and $p'=1$ initially, use
\begin{equation}
 S(u)=-\frac{c}{12}\log p'(u),\qquad
 F(u)=-\frac{c}{24\pi}\{p,u\}
 =\frac1{2\pi}\left(S''+\frac6c S'^2\right).
 \label{eq:BS}
\end{equation}
Here $u$ is an affinely normalized outgoing null coordinate, primes denote
$u$ derivatives, and $S$ is vacuum-subtracted entropy, not an absolute entropy.
Defining
\begin{equation}
 K=\frac{c}{12\pi},\quad w=e^{6S/c},\quad R=e^{6S_{\max}/c}>1,
 \qquad F=K\frac{w''}{w},\quad p'=w^{-2},
 \label{eq:w}
\end{equation}
consider a complete excursion with $w=1$ outside $[0,T]$, $1\le w\le R$,
and $\max w=R$. Nonnegative entropy relative to the endpoint value is an
explicit assumption. Allow $w\in W^{2,\infty}(\R)$, so $w,w'$ are continuous
and the flux is essentially bounded. Impose
\begin{equation}
 -\mathcal F_-\le F\le\mathcal F_+,\qquad
 a=\sqrt{\mathcal F_+/K}>0,\quad b=\sqrt{\mathcal F_-/K}>0.
 \label{eq:caps}
\end{equation}
These are externally specified caps, not pointwise energy conditions of QFT.

\section{Exact duration bound}

\begin{proposition}
Every excursion satisfying the above assumptions obeys
\begin{equation}
 \boxed{\ T\ge\mathcal T(a,b;R)
 =2\left[\frac{\arcosh y}{a}+\frac{\arccos(y/R)}{b}\right],\qquad
 y^2=\frac{b^2R^2+a^2}{a^2+b^2}.\ }
 \label{eq:bound}
\end{equation}
The bound is attained in the stated weak class and is the infimum over smooth
conformal-vacuum excursions with the same height and caps. Extra turnovers
cannot lower it.
\end{proposition}

\begin{proof}
Equation \eqref{eq:caps} is equivalent to $-b^2w\le w''\le a^2w$.
On a connected component $(\alpha,\beta)$ of $\{w'>0\}$, the endpoint
velocities vanish. Put $m=w(\alpha)\ge1$ and $M=w(\beta)\le R$.
Integrating the two acceleration inequalities, respectively from $\alpha$
and backwards from $\beta$, gives
\begin{equation}
 w'^2\le a^2(w^2-m^2),\qquad
 w'^2\le b^2(M^2-w^2).
\end{equation}
Time reversal gives the same estimates on decreasing components. Thus,
everywhere by continuity,
\begin{equation}
 |w'|\le V(w),\qquad
 V(x)=\min\{a\sqrt{x^2-1},\ b\sqrt{R^2-x^2}\}.
 \label{eq:envelope}
\end{equation}
Set $H(x)=\int_1^x ds/V(s)$. Its endpoint singularities are integrable.
The composition $H(w(u))$ is 1-Lipschitz: this follows first with $V$ replaced
by $\max(V,\epsilon)$, using the chain rule, and then by taking
$\epsilon\downarrow0$. If $w(u_*)=R$, then $H(R)\le u_*$ and
$H(R)\le T-u_*$. Consequently $T\ge2H(R)$, regardless of the number of
stationary points. The branches of \eqref{eq:envelope} cross at $y$, and
\begin{equation}
 H(R)=\int_1^y\frac{dx}{a\sqrt{x^2-1}}
      +\int_y^R\frac{dx}{b\sqrt{R^2-x^2}}
 =\frac{\arcosh y}{a}+\frac{\arccos(y/R)}b.
\end{equation}
This proves \eqref{eq:bound}.
\end{proof}

For attainment, write $t_+=\arcosh(y)/a$, $t_-=\arccos(y/R)/b$,
$\tau=t_++t_-$, and $T_*=2\tau$. The complete optimizer is
\begin{equation}
 w_*(u)=
 \begin{cases}
  1,&u\notin[0,T_*],\\
  \cosh(au),&0\le u\le t_+,\\
  R\cos[b(u-\tau)],&t_+\le u\le T_*-t_+,\\
  \cosh[a(T_*-u)],&T_*-t_+\le u\le T_*.
 \end{cases}
 \label{eq:optimizer}
\end{equation}
The identity $a\sqrt{y^2-1}=b\sqrt{R^2-y^2}$ ensures continuity of $w_*'$.
Its flux is $+\mathcal F_+$ on both outer intervals and $-\mathcal F_-$
on the central interval of length $2t_-$. Jumps of $w_*''$ produce finite
flux steps, not impulses. Equality in the Lipschitz argument forces
$H(w(u))=\min(u,T_*-u)$; hence the optimizer is unique on its support.
The net emitted energy satisfies
\begin{equation}
 E_{\rm net}=2\mathcal F_+t_+-2\mathcal F_-t_-
 =K\int_0^{T_*}(w_*'/w_*)^2\,du>0.
 \label{eq:energy}
\end{equation}

Smooth sharpness needs care because smoothing lowers the peak. Let
$z_\delta=\rho_\delta*w_*$, with a nonnegative smooth mollifier supported
in $[-\delta,\delta]$, and set
\begin{equation}
 \lambda_\delta=\frac{R-1}{\max z_\delta-1},\qquad
 w_\delta(u)=1+\lambda_\delta
 \left[z_\delta\left(u/\sqrt{\lambda_\delta}-\delta\right)-1\right].
 \label{eq:smooth}
\end{equation}
Convolution preserves $-b^2z_\delta\le z_\delta''\le a^2z_\delta$.
Here $\lambda_\delta\ge1$, $w_\delta\ge z_\delta$ at corresponding
arguments, and $w_\delta''=z_\delta''$. Both original caps therefore hold
exactly, while $\max w_\delta=R$. Its duration is
$\sqrt{\lambda_\delta}(T_*+2\delta)\to T_*$. Each $p_\delta'=w_\delta^{-2}$
defines a smooth, timelike moving mirror, asymptotically at rest. The final
affine intercept of $p_\delta$ may differ from the initial one. No globally
$C^2$ excursion attains $T_*$ because \eqref{eq:optimizer} has jumps of $w''$;
any duration strictly above $T_*$ is achievable by smoothing and dilation.

\section{Peak flux and the size of an excursion}

For equal caps $\mathcal F_+=\mathcal F_-=\mathcal F$, put $x=6S_{\max}/c$.
Equation \eqref{eq:bound} becomes the sharp peak-flux inequality
\begin{equation}
 \boxed{\ \mathcal F\ge\frac{c}{3\pi T^2}h(x)^2,\qquad
 h(x)=\frac12\arcosh(e^{2x})+\arctan\sqrt{\tanh x}.\ }
 \label{eq:peak}
\end{equation}
The two terms are the acceleration and braking times of the quantum-piston
solution \cite{Stefanatos}, after translating piston expansion into
exponentiated conformal entropy. In particular,
\begin{align}
 h(x)&=2\sqrt{x}+\tfrac1{15}x^{5/2}+O(x^{9/2}),&&x\downarrow0,\\
 h(x)&=x+\tfrac12\log2+\tfrac\pi4+O(e^{-2x}),&&x\to\infty.
\end{align}
Thus small excursions require
$\mathcal F_{\min}\sim8S_{\max}/(\pi T^2)$, whereas large excursions have
$\mathcal F_{\min}\sim12S_{\max}^2/(\pi cT^2)$.
The change from linear to quadratic dependence distinguishes a finite peak
constraint from merely bounding total emitted energy.

As a nonoptimal comparison, take $S=(c/6)x\sin^2(\pi u/T)$ on $[0,T]$,
with constant continuation. Its flux is
\begin{equation}
 \frac{F}{K}=\frac{\pi^2}{T^2}
 \left[2x\cos(2\pi u/T)+x^2\sin^2(2\pi u/T)\right].
\end{equation}
Its largest absolute flux, divided by the optimum \eqref{eq:peak}, is
\begin{equation}
 \frac{\mathcal F_{\sin^2}}{\mathcal F_{\min}}
 =\frac{\pi^2}{4h(x)^2}
 \begin{cases}2x,&0<x\le1,\\x^2+1,&x\ge1.\end{cases}
 \label{eq:sin}
\end{equation}
The ratio tends to $\pi^2/8$ for small excursions and $\pi^2/4$ for large
ones. This comparison concerns peak flux; it is different from the
integrated-energy sinusoidal comparison in Ref.~\cite{KD}.

Unequal caps expose two complementary obstructions:
\begin{equation}
 \lim_{b\to\infty}\mathcal T=\frac{2\arcosh R}{a},\qquad
 \lim_{a\to\infty}\mathcal T=\frac{2\arccos(1/R)}{b}.
 \label{eq:limits}
\end{equation}
At fixed positive cap, vanishing negative cap sends the duration to infinity.
At fixed negative cap, allowing arbitrarily intense positive pulses leaves
a nonzero braking time. The one-sided limits admit impulses of the unconstrained
sign and should be interpreted as limits, not finite-cap optimizers.

\section{A diagnostic for QNEC nonsaturation}

For a state and null deformation obeying the strengthened two-dimensional
quantum null energy condition (QNEC), write
\begin{equation}
 F=F_S+\Delta_Q,\qquad
 F_S=\frac1{2\pi}\left(S''+\frac6c S'^2\right),\qquad
 0\le\Delta_Q\le D.
 \label{eq:gap}
\end{equation}
Here the condition must hold for the same entropy $S$ used above, for example
in a half-line limit where the reference entropy has vanishing deformation
derivatives. Generic finite-interval vacuum subtraction cannot be inserted
into this nonlinear condition without justification.
Saturation and nonsaturation of this condition distinguish physically different
classes of two-dimensional states \cite{Ecker}. With the same entropy
boundary conditions and measured flux caps, \eqref{eq:gap} implies
\begin{equation}
 -(\mathcal F_-+D)\le F_S\le\mathcal F_+.
\end{equation}
The mathematical proposition therefore gives the necessary condition
\begin{equation}
 T\ge\mathcal T\left(\sqrt{\mathcal F_+/K},
                     \sqrt{(\mathcal F_-+D)/K};R\right).
 \label{eq:gapbound}
\end{equation}
For fixed $a,R$, $\mathcal T(a,b;R)$ is strictly decreasing in $b$:
in \eqref{eq:envelope} the braking envelope increases on a nonempty interval.
If
$2\arcosh(R)/a<T<\mathcal T(a,\sqrt{\mathcal F_-/K};R)$,
there is a unique $b_{\rm crit}>\sqrt{\mathcal F_-/K}$ solving
$\mathcal T(a,b_{\rm crit};R)=T$. Consequently
\begin{equation}
 \esssup\Delta_Q\ \ge\ Kb_{\rm crit}^{2}-\mathcal F_->0.
 \label{eq:witness}
\end{equation}
An excursion at or below the positive-cap floor $2\arcosh(R)/a$ is incompatible
with a regular profile under these assumptions: equality would require an
instantaneous reversal of nonzero $w'$ at the maximum. For durations at or above the zero-gap bound this
diagnostic supplies no positive lower bound. Equation \eqref{eq:witness} is
a necessary witness, not a claim that every pair $(S,F)$ obeying QNEC is
realized by a quantum state, nor a claim of sharpness in the nonsaturating sector.
For a smooth profile a positive threshold in \eqref{eq:witness} is a strict
inequality, since equality would force \eqref{eq:optimizer}.

The diagnostic becomes strong near the positive-cap floor
$T_{\rm f}=2\arcosh(R)/a$. Expanding \eqref{eq:bound} at large $b$ gives
\begin{equation}
 \mathcal T=T_{\rm f}+\frac{a\sqrt{1-R^{-2}}}{b^2}+O(b^{-4}),\qquad
 D_{\rm crit}\sim\frac{Ka\sqrt{1-R^{-2}}}{T-T_{\rm f}}.
 \label{eq:divergence}
\end{equation}
For example, $R=e$ and $\mathcal F_+=\mathcal F_-=K$ give
$T_{\rm f}=3.314909$ and a zero-gap duration $4.123581$. An excursion of
duration $3.5$ requires $D\ge3.853457K$. This is a conditional diagnostic,
not an example of a constructed nonsaturating state.

\section{Scope and interpretation}

The result supplies a benchmark for conformal-vacuum Page-like pulses with
finite stress: it fixes the optimal duration, the switch locations, and the
minimum peak magnitude. Its possible use beyond that sector is the conditional
QNEC-gap diagnostic. The formulas refer to the same affine null coordinate
throughout. Under $\widetilde u=\lambda u$,
$\widetilde T=\lambda T$ and $\widetilde{\mathcal F}_{\pm}=
\mathcal F_{\pm}/\lambda^2$, so the bound is covariant under this rescaling.

The flux caps are assumptions on expectation values and do not follow from
quantum energy inequalities, which constrain smeared energy \cite{FH}.
Neither entropy undershoots below the endpoint value nor thermal backgrounds
are covered. The constructed mirror is a prescribed boundary, not a solution
of gravitational backreaction. In particular, an island formula for black-hole
entropy cannot simply be substituted into \eqref{eq:BS}. These restrictions
leave a concrete CFT optimization and diagnostic; they do not establish a
universal evaporation-time bound.

\begin{thebibliography}{9}
\bibitem{BS} E.~Bianchi and M.~Smerlak,
``Entanglement entropy and negative energy in two dimensions,''
Phys. Rev. D \textbf{90}, 041904 (2014),
\href{https://arxiv.org/abs/1404.0602}{arXiv:1404.0602}.
\bibitem{KD} S.~Karmoudi and M.~Daoud,
``Sharp null-energy lower bounds for Page-like entropy excursions in
two-dimensional conformal field theory,''
\href{https://arxiv.org/abs/2610.03576v2}{arXiv:2610.03576v2} (2026).
\bibitem{AP} S.~Abdolrahimi and D.~N.~Page,
``Hawking radiation energy and entropy from a Bianchi--Smerlak semiclassical
black hole,'' Phys. Rev. D \textbf{92}, 083005 (2015),
\href{https://arxiv.org/abs/1506.01018}{arXiv:1506.01018}.
\bibitem{Stefanatos} D.~Stefanatos,
``Optimal shortcuts to adiabaticity for a quantum piston,''
Automatica \textbf{49}, 3079--3083 (2013),
\href{https://arxiv.org/abs/1301.5137}{arXiv:1301.5137}.
\bibitem{TW} E.~Teo and K.~F.~Wong,
``Quantum interest in two dimensions,'' Phys. Rev. D \textbf{66}, 064007 (2002),
\href{https://arxiv.org/abs/gr-qc/0206066}{arXiv:gr-qc/0206066}.
\bibitem{Ecker} C.~Ecker, D.~Grumiller, W.~van der Schee,
M.~M.~Sheikh-Jabbari, and P.~Stanzer,
``Quantum null energy condition and its (non)saturation in 2d CFTs,''
SciPost Phys. \textbf{6}, 036 (2019),
\href{https://arxiv.org/abs/1901.04499}{arXiv:1901.04499}.
\bibitem{FH} C.~J.~Fewster and S.~Hollands,
``Quantum energy inequalities in two-dimensional conformal field theory,''
Rev. Math. Phys. \textbf{17}, 577--612 (2005),
\href{https://arxiv.org/abs/math-ph/0412028}{arXiv:math-ph/0412028}.
\end{thebibliography}
\end{document}
